\section{The New Prophecy and the Cost of Rigor}

\subsection{No single conversion in 207}

Older biographies often divide the corpus into Catholic works before 207 and
Montanist works afterward, then supply a conversion, schism, or excommunication
at the dividing line. The books do not preserve that event. They show an
increasingly explicit commitment to the prophetic movement associated with
Montanus, Prisca, and Maximilla; they do not give the day Tertullian joined a
separate institution, the act by which a bishop expelled him, or one uncontested
chronology of every work.

``New Prophecy'' is the most responsible historical name for the movement's own
claim. ``Montanism'' is conventional but can suggest that one founder, doctrine,
and organization explain every community across Phrygia, Rome, and Carthage.
Tertullian's African reception of prophetic sayings and disciplines need not
reproduce the movement's first Asian form. Opponents' reports of ecstasy,
eschatology, money, women, and failed prophecy must be read beside adherents'
fragments rather than fused into one hostile portrait.

Eusebius preserves the principal early anti-New-Prophecy dossier in
\emph{Ecclesiastical History} 5.16--19. He excerpts an anonymous disputant,
Apollonius, Serapion, and related letters from controversies nearer the
movement's Asian setting. That proximity makes the dossier indispensable; its
hostile selection and Eusebius's later excerpting make it mediated evidence.
Reports of ecstatic speech, money, deaths, and failed prediction are not
unfiltered minutes, and none narrates Tertullian's African adherence.

The trajectory is thematic before it is calendrical. Early writings already
value martyrdom, sexual restraint, fasting, public repentance, and visible
separation from idols. Later writings invoke the Paraclete to intensify these
practices, reject second marriage, extend fasts, enforce virginal veiling, and
deny ecclesial absolution for grave sins. Prophecy did not create rigor from
nothing. It gave development an authority and direction.

\subsection{Faith fixed, discipline advancing}

\emph{On the Veiling of Virgins} 1 states the architecture clearly. The rule of
faith is one, immovable, and irreformable; discipline and conduct may advance as
the Paraclete leads into truth. Tertullian illustrates growth through stages of
human life and divine education. The distinction allows him to deny doctrinal
innovation while demanding practices his opponents regard as innovations.

That logic deserves more than retrospective dismissal. It recognizes that a
living Church encounters questions not enumerated word for word in the initial
rule. It also poses the decisive problem: how is genuine development discerned,
and who judges a prophet's claim? Tertullian appeals to Scripture, apostolic
faith, moral fruits, vision, and the promised Spirit, but his increasingly
adversarial language weakens the possibility of communal correction. The
Paraclete can become the warrant for the strictest available conclusion.

In \emph{Against Praxeas} 1, Tertullian says a bishop of Rome had recognized the
prophecies of Montanus, Prisca, and Maximilla and was bringing peace to the
churches of Asia and Phrygia until Praxeas induced him to recall his letters.
No pope is named. Attempts to identify Victor, Zephyrinus, or another bishop
depend on disputed chronology. The story is partisan: it makes the theological
opponent responsible for both suppressing prophecy and confusing Father and
Son. It establishes Tertullian's account of a revoked Roman opening, not an
independently preserved papal dossier.

The same chapter says that recognition of the Paraclete later separated ``us''
from the \emph{psychici}. The Latin can mark withdrawal or division without
specifying a canonical sentence. This is the strongest autobiographical notice
of a breach, but it cannot supply the named judge, date, procedure, or stable
denominational map that it omits.

\subsection{Fasts, marriages, and a Church that may forgive}

\emph{On Fasting} defends additional fasts and dry-food observances against
Christians who regard them as novelty or excess. Tertullian treats appetite as
a school of obedience and biblical history as a pattern of restraint. His
opponents' concern---that a humanly intensified discipline may be imposed as
divine law---is not answered merely by the seriousness of fasting. The dispute
concerns authority as much as food.

The marriage sequence is still clearer. \emph{To His Wife} discourages a widow
from remarrying. \emph{On Exhortation to Chastity} acknowledges that a second
marriage may be permitted but presses the difference between lawful and
expedient. \emph{On Monogamy}, appealing to the Paraclete, rejects second
marriage as inconsistent with the final discipline. This is development in
Tertullian's own terms and hardening in historical description. The reader
should not collapse the three works into one timeless teaching.

The gravest contest concerns repentance. Early \emph{On Repentance} provides
for a difficult second repentance after baptism. Later \emph{On Modesty} 1
mocks an unnamed episcopal edict granting forgiveness for adultery and
fornication, calling its author the ``supreme pontiff'' and ``bishop of
bishops.'' Identification with Callistus of Rome is traditional but contested;
some reconstructions locate the target in Carthage. The satire therefore cannot
become a dated papal quotation without argument.

In chapter 21 Tertullian relocates the Church's decisive power from an
institution represented by a number of bishops to the Church of the Spirit,
acting through spiritual persons. He does not deny every ministry or community.
He denies that bishops as such may exercise the claimed absolution. The cost of
rigor is now ecclesiological: a baptismal fear of sin has become a principle by
which the visible Church's mercy proves its lack of spiritual authority.

\subsection{Prophecy, isolation, and the absence of an ending}

Tertullian's prophetic texts preserve phenomena not reducible to stricter rules.
\emph{On the Soul} 9 values a woman's visions examined after worship;
\emph{On the Resurrection of the Flesh} and other later works attribute
clarification to the Paraclete; \emph{Against Praxeas} presents the Spirit as a
teacher of Trinitarian distinction. For him the movement promised not escape
from catholic faith but its eschatological maturation.

Yet the rhetoric of spiritual against soulish Christians repeatedly divides
the baptized into those who can bear truth and those who compromise it. Mercy,
pastoral accommodation, and institutional discernment are readily redescribed
as appetite or cowardice. Tertullian's moral protest against easy Christianity
remains powerful; his confidence that maximal severity identifies the Spirit
does not.

Jerome later explains the break by envy and insults from the Roman clergy. That
story appears in \emph{On Illustrious Men} 53 about 180 years after Tertullian's
last securely datable work, without a named source. It may preserve lost
reception, but it cannot control motive. No extant work supplies a Roman
clerical career whose jealous colleagues drove a disappointed presbyter away.

Nor does the archive provide a final reconciliation or dramatic apostasy. The
last securely datable writing is \emph{To Scapula} after the eclipse of 212;
relative arguments place some polemics later, but no text closes the career.
The voice that had demanded an open hearing from governors disappears without
a trial record, deathbed, or last sentence. What remains is the breach inside
the books.
